\section{The mixed three-box ordinary alternating tail}
\label{sec:alternating-mixed-three-tail}

\paragraph{Status and scope.}
New root derivation HD-MT301--03. Complete author and peer
reviews are pending. The ordinary label is $(n-3,2,1)$,
with $n\ge15$. Put
\[
 a=\frac{n(n-2)(n-4)}3,\quad
 A=\frac{n^2-8n+8}3,\quad B=\frac{n(n-3)}3,\quad
 b=\frac n3,\quad d_0=A+B.                              \tag{MT3.1}
\]
Here $a$ is the integral degree of $W=S^{(n-3,2,1)}$;
the remaining constants are rational bounds. Its restriction
to $A_n$ is irreducible, as this long label is not self-conjugate.
The rational $S_n$ model and full compatible normalizer are
the accepted ordinary alternating interfaces.

The exact known character input is
\[
               \chi_W(x)=\frac{f(f-2)(f-4)}3-v,          \tag{MT3.2}
\]
where $f$ and $v$ count the one- and three-cycles of $x$.
It is the $(n-3,2,1)$ row of Garsia--Goupil,
\emph{Character Polynomials, their q-Analogs and the Kronecker
Product}, Table~5.1, p.~31 \cite{HD:GarsiaGoupil2009}.
The general character-polynomial statement is Proposition~I.1,
p.~2, whose complete invoked proof is Section~2, equations
(2.1)--(2.6), pp.~9--11. Those passages and the table were
freshly read. We cite this published formula instead of
reproving the known character theory.

\subsection*{HD-MT301. The all-eigenvalue index estimate}

\begin{lemma}[HD-MT301]
For every nonidentity permutation $x\in S_n$,
\[
 \nu(Wx)\ge A I(x)+B\ge d_0,
 \qquad \chi_W(x)\ge-b,                                \tag{MT3.3}
\]
where $I(x)=n-\#\operatorname{cycles}(x)$ includes fixed
letters, and $\nu$ is maximum-eigenspace codimension.
The character lower bound also holds at the identity.
\end{lemma}

\begin{proof}
The constraint $3v\le n-f$ gives
\[
 \chi_W(x)\ge\frac{f(f-2)(f-4)+f-n}3
            =\frac{f(f-3)^2-n}3\ge-\frac n3.             \tag{MT3.4}
\]
Also $v\ge0$, so for $f<n$,
\[
 \frac{a-\chi_W(x)}{n-f}\ge c_f,
 \qquad
 c_f=\frac{n^2+nf+f^2-6n-6f+8}3,
 \quad c=c_0=\frac{(n-2)(n-4)}3.                         \tag{MT3.5}
\]
This follows by exact polynomial division. Since $n\ge15$,
$c_f$ is increasing for $f\ge0$, and hence $c_f\ge c$.
At the identity the un-divided defect inequality has both
sides zero.

Fix $x$ of order $L$ and index $I$. Cyclic averaging gives
$a-m_1\ge cI$, where $m_1$ is its fixed-space multiplicity:
the average number of fixed letters equals its number of
cycles. The all-eigenvalue Fourier calculation (AT3.7),
using the just-proved real-character lower bound $-b$,
gives $m_z\le m_1+b$ for every root $z$. Explicitly,
\[
 m_z-m_1=\frac1L\sum_{j=0}^{L-1}\chi_W(x^j)
                 (\operatorname{Re}z^{-j}-1)\le b
 \quad(z\ne1),                                         \tag{MT3.6}
\]
and the inequality is immediate for $z=1$.
Thus $\nu(Wx)\ge cI-b$.

If $I\ge(n-2)/2$, then $c-A=2n/3$ implies
\[
 \nu(Wx)-AI\ge\frac{2nI-n}{3}
                         \ge\frac{n(n-3)}3=B.           \tag{MT3.7}
\]
Otherwise every power of $x$ fixes at least $f_0=n-2I>2$
letters, since the moved support of a permutation is at
most twice its index. Monotonicity of $c_f$ improves the
cyclic bound to $a-m_1\ge c_{f_0}I$.
After using (MT3.6) and clearing positive denominators,
$c_{f_0}I-b\ge AI+B$ is exactly
\[
 (n-f_0)\bigl(f_0^2+(n-6)f_0+2n\bigr)\ge2n(n-2).
                                                                  \tag{MT3.8}
\]
The difference of the two sides is
\[
 p(f_0)=-f_0^3+6f_0^2+(n^2-8n)f_0+4n.                   \tag{MT3.9}
\]
It is concave on $[2,n-2]$, because $p''(f)=12-6f\le0$.
Its endpoint values are $2(n-2)(n-4)$ and $2(n-4)^2$,
both strictly positive. It is therefore positive throughout
that interval. This proves the required uniform bound for
every nonidentity permutation and every target eigenvalue.
\end{proof}

\subsection*{HD-MT302. The one-factor source in every field}

\begin{theorem}[HD-MT302]
The actual rank-nine sixfold inducing source cannot have
a single ordinary core $A_n$ on $S^{(n-3,2,1)}$, $n\ge15$.
This permits all compatible scalar and outer actions over
every original imaginary quadratic field.
\end{theorem}

\begin{proof}
The rational $S_n$ normalizer meets the scalars trivially:
a scalar centralizes the faithful $A_n$ image, while
the centralizer of $A_n$ in $S_n$ is trivial. Extracting
the natural permutation is therefore an actual homomorphism,
and its image contains $A_n$. The associated quotient
cover is connected and has degree $n$.

The actual Euler identity is $6=2a(h-1)+D$.
Since $a\ge715$, genus $h\ge2$ is impossible, and
$h=1$ permits at most six genuine marks and period rank
at most six. Thus rank nine requires $h=0$, at least
twelve genuine inertia entries, and $D=2a+6$.
Riemann--Hurwitz for the actual degree-$n$ quotient gives
$\sum I(x_i)\ge2n-2$.
Each nonscalar entry moves at least $AI(x_i)+B$ dimensions
by (MT3.3). A genuine scalar moves all $a\ge B$ dimensions
and has $I=0$, so it satisfies the same bound.
Consequently
\[
 D-2a\ge2A(n-1)+12B-2a
       =\frac{6n^2-20n-16}{3}>6,                         \tag{MT3.10}
\]
contradicting the source Euler identity. No scalar-conductor
assumption is used in this one-factor argument.
\end{proof}

\subsection*{HD-MT303. All transitive tensor counts}

\begin{theorem}[HD-MT303]
Let the actual inducing image normalize the ordinary core
$A_n^t$ on $W^{\otimes t}$ transitively on its factors,
where $W=S^{(n-3,2,1)}$, $n\ge15$ and $t\ge2$.
If its character field contains a non-Eisenstein imaginary
quadratic field, it cannot give the rank-nine sixfold source.
All compatible outer automorphisms, scalar phases and
factor permutations are included.
\end{theorem}

\begin{proof}
Write $e=a^t$. The faithful rational tensor-wreath normalizer
has trivial scalar intersection, by its conjugation action
on the individual faithful simple factors. Thus its scalar
and rational coordinates are actual homomorphisms, and its
natural $nt$-letter action is transitive. The inducing
character field is contained in the cyclotomic field of
the scalar character. The original non-Eisenstein field
therefore forces at least two inertia phases outside $\mu_6$,
by generation and product one. It need only be contained
in the inducing field; equality is never assumed.

The actual Euler and marked-cover arguments again give
$h=0$, $r\ge12$, $D=2e+6$.
Every power of a rational tensor-wreath operator has trace
at least $-(b/a)e$, since its trace is a product of
cycle-product traces bounded below by $-b$ and in absolute
value by $a$. Rational cyclotomic multiplicities and the
fourth-root square-trace calculation of D-ATT-02 give
moving rank at least $(3-b/a)e/4$ at both scalar marks.
Every genuine inertia has support at least $(d_0/a)e$:
use (MT3.3) and spectral convolution if factors are fixed;
if they move, the weighted-word argument gives the stronger
$(a-1)e/(2a)$. Directly from (MT3.1),
$2d_0\le a-1$, $b\le d_0$, and $d_0>97$ for $n\ge15$.

If three entries move factors, choose one distinct from
both marked scalar entries. With nine further genuine marks,
\[
 D\ge2e+a^{t-1}\left(9d_0-\frac{b+1}{2}\right)>2e+6.
                                                                  \tag{MT3.11}
\]
Therefore exactly two entries move factors: zero contradicts
transitivity, and one contradicts product one. Their factor
permutations are inverse full $t$-cycles.

Let $H$ be the sum of the natural indices of their two
cycle products and $N$ the sum of all natural indices
at the factor-fixing entries. Riemann--Hurwitz for the
actual degree-$nt$ quotient gives
$2n(t-1)+H+N\ge2nt-2$, hence $H+N\ge2n-2$.
At a factor-fixing genuine entry, averaging the active
factor estimates in (MT3.3) gives support at least
$a^{t-1}(AI/t+B)$, preserving the full $B$.
A genuine scalar has full rank and satisfies this with
$I=0$. At each full factor cycle the exact D-ATT-08 bound
is
\[
 (1-1/t)e+\frac{a^{t-1}}t\nu(Wx)-\frac{R_t(a)}t.
                                                                  \tag{MT3.12}
\]
Here use only $\nu(Wx)\ge AI(x)$, which is also valid
for the identity cycle product. The exact errors satisfy
$R_2(a)=a$ and $R_t(a)\le2a^{t-2}$ for $t\ge3$.
Thus $2R_t(a)/a^{t-1}\le2$ for every $t\ge2$.
There are at least ten factor-fixing entries. Summing gives
\[
 \begin{split}
 D-2e
 &\ge\frac{a^{t-1}}t
 \left(2A(n-1)-2a+10tB-2\right)\\
 &=\frac{a^{t-1}}t
 \left(10tB-\frac{6n^2-16n+16}{3}-2\right)
 >B a^{t-1}>6.
 \end{split}                                             \tag{MT3.13}
\]
Indeed $9tB\ge18B$ and
\[
 18B-\frac{6n^2-16n+16}{3}-2
                    =\frac{12n^2-38n-22}{3}>0
 \quad(n\ge15).                                         \tag{MT3.14}
\]
This contradicts the Euler budget and proves the theorem.
\end{proof}

\paragraph{Strictness and interpretation.}
The earlier exact shape gate allows $(97,2,1)$ at $n=100$
over $\mathbf Q(i)$: its tail has size three and its
ALT07 fraction is less than $1/10$ even in the tensor limit.
This label is distinct from the already excluded standard,
two-box and two-row three-box labels. The new result
removes another entire necessary supplier family, without
asserting that an actual source previously existed there.
It constrains actual finite-cover local systems and their
representation lambda-rings and K-theory; in the project's
frame it shrinks condition-A suppliers (cote16scan050;
SIGN_NEXT SN-C03/C04). No condition-I cycle or new conditional
Markman, Severn or032 input is produced.
